\documentclass[profile=standard]{ipara-handbook}
\title{B02｜Jacobian 与行列式：坐标变换与局部体积缩放}
\author{}
\date{}
\begin{document}
\maketitle

\begin{quote}
本册回答：\textbf{为什么一个非线性坐标变换最后会在积分微元前出现一个行列式？为什么有时乘 Jacobian、有时乘逆 Jacobian，而这件事不应该靠背“倒数公式”解决？}
\end{quote}

\section{本册定位：把局部线性作用压缩成体积倍率}
B01 Own “可微产生局部线性映射”；高数 Topic07 Own Jacobian matrix；线代 Topic02 Own determinant 的代数与几何机制；高数 Topic08 Own 多重积分换元。B02 只 Own 它们之间的接口：
\[
\boxed{
T\text{ 可微}
\xrightarrow{\text{B01}}
DT_u
\xrightarrow{\text{选基}}
J_T(u)
\xrightarrow{\det}
\text{局部有向体积倍率}
\xrightarrow{\lvert\cdot\rvert}
\text{无向面积/体积微元倍率}
}
\]

如果没有这座桥，Jacobian 很容易退化成一个“换元时别忘了乘的东西”，学生也就无法自行判断变换方向、奇异点和多分支。

\section{先分清两个对象：Jacobian matrix 与 determinant}
对可微映射
\[
T:\mathbb R^n\to\mathbb R^m,
\]
Jacobian matrix $J_T(u)$ 是导数线性映射 $DT_u$ 在所选坐标基下的矩阵表示，维度为 $m\times n$。它记录每个输入方向局部被送到哪里。

只有 $m=n$ 时，才可以继续取 Jacobian determinant
\[
J(u):=\det J_T(u).
\]
这里先固定符号边界：本册写 $J_T$ 时指完整 Jacobian matrix，写 $\det J_T$ 或 $J(u)$ 时才指 determinant。行列式把一张完整的线性映射压缩成一个标量：局部满维小平行体的\textbf{有向体积倍率}；其绝对值才是普通无向面积/体积倍率。

因此：
\[
\boxed{J_T\neq\det J_T.}
\]
前者保留方向作用，后者只保留满维缩放与定向信息。

\section{为什么 determinant 正好是微元倍率}
在二维中，一个很小的 $du\,dv$ 矩形由边向量
\[
e_1\,du,\qquad e_2\,dv
\]
张成。经过局部线性映射 $DT_u$ 后，两条边近似变成
\[
DT_u(e_1)\,du,\qquad DT_u(e_2)\,dv.
\]
它们张成的小平行四边形有向面积是原面积的
\[
\det J_T(u)
\]
倍。因此普通面积只取大小：
\[
\boxed{\mathrm{d}A_x=\lvert\det J_T(u)\rvert\,\mathrm{d}A_u.}
\]
三维完全同理：determinant 是局部体积倍率。这里的微元式是局部缩放的压缩写法；真正的换元定理还必须同时控制区域覆盖与映射正则性。

\section{换元公式不是“公式搬家”，而是总累积量保持不变}
若 $(x,y)=T(u,v)$，原积分为
\[
\iint_{D_{xy}} f(x,y)\,\mathrm{d}x\,\mathrm{d}y.
\]
换坐标以后，同一块几何区域被另一套网格编码。若 $T$ 在所用分片上足够光滑且一一，则
\[
\boxed{
\iint_{D_{xy}}f(x,y)\,\mathrm{d}x\,\mathrm{d}y
=
\iint_{D_{uv}}f(T(u,v))\,
\lvert\det J_T(u,v)\rvert\,\mathrm{d}u\,\mathrm{d}v.
}
\]
这里真正保持的是\textbf{累积的几何/物理总量}；变的是区域编码、函数表示和局部微元。

\section{为什么正变换与逆变换不能靠记忆选倒数}
假设题目先给
\[
(u,v)=S(x,y),
\]
但积分最终要写在 $(u,v)$ 中。最稳的办法不是问“Jacobian 要不要取倒数”，而是先写你实际使用的映射方向：
\[
(x,y)=T(u,v)=S^{-1}(u,v).
\]
那么直接使用
\[
\left\lvert\frac{\partial(x,y)}{\partial(u,v)}\right\rvert.
\]
若两边局部互逆，则链式法则给
\[
J_TJ_S=I,
\]
因此
\[
\det J_T=\frac{1}{\det J_S}.
\]
“倒数”只是逆映射线性化后的结果，不是另一条孤立规则。

\section{组合变换：为什么倍率会相乘}
若先做 $T$ 再做 $S$，B01 给出
\[
J_{S\circ T}=J_S(T(u))J_T(u).
\]
取 determinant：
\[
\det J_{S\circ T}
=\det J_S(T(u))\det J_T(u).
\]
因此连续两次坐标缩放的局部体积倍率相乘。这也提供了复杂换元的独立校验：若把一个变换拆成两步，最终倍率应与直接求复合 Jacobian 一致。

\section{母例：极坐标中的 $r$ 从哪里来}
令
\[
T(r,\theta)=(r\cos\theta,r\sin\theta).
\]
则
\[
J_T=
\begin{pmatrix}
\cos\theta&-r\sin\theta\\
\sin\theta&r\cos\theta
\end{pmatrix},
\qquad
\det J_T=r.
\]
于是
\[
\mathrm{d}A=r\,\mathrm{d}r\,\mathrm{d}\theta.
\]
这个 $r$ 不是人为补偿：角度变化 $\mathrm{d}\theta$ 在半径 $r$ 处对应的切向长度本来就是约 $r\,\mathrm{d}\theta$，与径向厚度 $\mathrm{d}r$ 相乘自然形成面积 $r\,\mathrm{d}r\,\mathrm{d}\theta$。

在 $r=0$ 处 determinant 为零，说明角度坐标失去区分能力：所有 $\theta$ 都描述同一个原点。这正是“坐标退化”的几何含义。注意这里的零点本身并不让整个极坐标积分失效；它是一个需要由换元定理和区域结构共同处理的退化点。

\section{局部资格不等于全局覆盖}
\subsection{$\det J_T\neq0$ 只给局部可逆性线索}
这里 $C^1$ 指一阶导数连续，也就是一阶连续可微。
若 $T$ 在点附近为 $C^1$，且该点满足 $\det J_T\neq0$，逆函数定理给出局部可逆。它说明局部没有把满维方向压扁，但\textbf{不能单独保证整个区域上一一}。

例如极坐标若让 $\theta$ 横跨多个完整周期，绝大多数点会被重复覆盖，即使 $r>0$ 处 Jacobian 非零。

\subsection{多对一必须分片或计重}
普通换元公式默认在所使用区域上控制好覆盖次数。若一个新坐标点对应多个原坐标分支，应拆区或显式累计各分支，不能只算一遍 determinant。

\subsection{determinant 为零意味着局部降维}
若 $\det J_T=0$，局部小体积被压到低维，普通满维换元不能机械通过这一点。题目若只在零测集上出现孤立奇异点，有时整体积分仍可处理；但必须由具体定理/区域决定，不能把“某一点为零”与“整个换元必失败”简单等同。

\section{B02 串起了哪些内容}
\begin{tabular}{p{0.23\linewidth}p{0.31\linewidth}p{0.38\linewidth}}
\hline
Owner / 下游 & 对象 & B02 的接口\\
\hline
B01 / 高数 Topic07 & 局部线性映射、Jacobian matrix & 只在同维方阵时取 determinant 读取满维缩放\\
线代 Topic02 & determinant & 将“线性映射体积倍率”翻译到非线性坐标变换的局部网格\\
高数 Topic08 & 重积分区域与微元 & 输出新坐标微元倍率，并要求同步搬运区域\\
B07 / 概率 Topic04 & 随机变量变换 & 把同一体积倍率进一步解释为密度必须反向补偿的概率质量守恒\\
\hline
\end{tabular}

\section{反桥梁：四个高频偷换}
\begin{tabular}{p{0.20\linewidth}p{0.04\linewidth}p{0.20\linewidth}p{0.48\linewidth}}
\hline
A & & B & 真正区别\\
\hline
Jacobian matrix & $\neq$ & determinant & 前者可矩形且保留方向作用；后者仅方阵且压成体积标量\\
$\det J\neq0$ & $\neq$ & 全局一一 & 非零只控制局部退化，覆盖仍需全局检查\\
正变换 Jacobian & $\neq$ & 逆变换 Jacobian & 二者对应不同映射方向，互为倒数需先有局部可逆\\
有向体积 & $\neq$ & 普通面积/体积 & determinant 带方向符号；积分微元通常取绝对值\\
\hline
\end{tabular}

\section{调用信号与第一观察}
题面出现“二重/三重积分换元、极坐标/柱坐标、给出新旧变量关系、变量变换的密度”时，先写四行账本：
\begin{enumerate}
\item 当前实际使用的映射方向：旧变量 = 新变量的函数，还是相反？
\item 新区域是不是原区域的完整逆像/像，是否一一覆盖？
\item Jacobian matrix 是否方阵，determinant 是否退化？
\item 当前任务只需几何微元，还是要继续转 B07 维护概率质量？
\end{enumerate}

\section{一页压缩}
B02 的生成句只有一句：
\[
\boxed{\text{非线性换坐标在足够小尺度上由线性映射控制，而 determinant 正是这张线性映射的满维体积倍率。}}
\]
因此忘掉所有换元口诀时，只需从
\[
dV_{\rm old}\approx\lvert\det DT\rvert\,dV_{\rm new}
\]
重建公式，再同时核对变换方向、区域覆盖、绝对值与奇异点。

\end{document}
